% Calcul des quantités de matière lors d'une réaction
À $\qty{600}{\degreeCelsius}$, le carbone réagit avec l'oxyde de cuivre(II)
\ch{CuO} pour former du dioxyde de carbone et du cuivre. On met en présence
$\qty{1.00}{\g}$ d'oxyde de cuivre et $\qty{0.120}{\g}$ de carbone.
\begin{enumerate}
    \item Quel est le réactif limitant~?
    \item Quel est l'avancement maximal~?
    \item Quelles sont les quantités de matière de dioxyde de carbone et de
          cuivre formés~?
\end{enumerate}
\begin{donnees}
    $M_{\ch{C}}=\qty{12.0}{\g\per\mol}$~;
    $M_{\ch{Cu}}=\qty{63.5}{\g\per\mol}$~;
    $M_{\ch{O}}=\qty{16.0}{\g\per\mol}$.
\end{donnees}


